\printedpage{49}
% head: A. Grothendieck : Dérivateurs, Ch. II, Édition des Universités de Montpellier II et Paris VII
% foot: Version 8 juillet 2001

\begin{itemize}[leftmargin=2.5em]
\item[5)] Let $K$ be a field (for simplicity), $L$ a Galois extension with group $G$. Relation between $L$ and $G$
\[
\gamma\cdot x=x,\qquad \text{i.e. } x\in L^{\gamma}.
\]
The closed subsets of $L$ are the subextensions, the closed subsets of $G$ the subgroups of $G$.
\end{itemize}

There are various analogues in inseparable cases (by dint of $p$-Lie algebras for example), and descent statements of all kinds.

\msp{page 80}All these examples are of the 1\textsuperscript{st} kind, the ordered sets that occur are of the form $\Pf(P_0)$, $\Pf(Q_0)$, and the relation between them is deduced from a relation $R_0$ between $P_0$ and $Q_0$. Here now are more complex examples.

\begin{itemize}[leftmargin=2.5em]
\item[6)] Let $\mathcal{M}$ be a category (``big'', most often \dots). We take
\[
\mathcal{P}=\mathcal{Q}=\Pf(\Fl(\mathcal{M})).
\]
The relation $\mathcal{R}\subseteq\mathcal{P}\times\mathcal{Q}$ is Quillen's relation, which I write, for $\Phi,\Psi\in\Pf$
\[
\Psi\subseteq\Phi_{*},
\]
or again equivalently
\[
\Phi\subseteq\Psi^{*},
\]
($\Psi$ has the LLP \texttt{[left lifting property]} with respect to $\Psi$, $\Phi$ the RLP \texttt{[right lifting property]} with respect to $\Psi$). Thus the maps $\alpha_0$, $\beta_0$ are given by % [sic] Psi a la LLP pour Psi -> Psi a la LLP pour Phi, the relation is between Psi and Phi (cf. "Phi la RLP pour Psi")
\[
\beta_0(\Phi)=\Phi_{*},\qquad \alpha_0(\Psi)=\Psi^{*}.
\]
Here it would occur to no one to introduce $\Pf(\mathcal{P})$, $\Pf(\mathcal{Q})$ - it is already big enough as it is! The notation for $\Phi^{\natural}$ is here $\tilde{\Phi}$, for $\Psi^{\flat}$ it is $\bar{\Psi}$. The closed subsets are those I call Q-closed on the left and on the right.

The relation $\mathcal{R}$, which I write
\[
\Phi\longleftrightarrow\Psi,
\]
\msp{page 81}is not deduced from a relation between arrows of $\mathcal{M}$, i.e. from a subset of $\Fl(\mathcal{M})\times\Fl(\mathcal{M})$, but is of a more complex nature. To describe it, one must introduce the set $\mathcal{D}$ of commutative square diagrams (solid lines) of the form
\[
\begin{tikzcd}
A \arrow[r] \arrow[d, "i"] & X \arrow[d, "p"] \\
B \arrow[r] \arrow[ur, dotted, "f"] & Y
\end{tikzcd}
\tag*{$(*)$}
\]
\printedpage{50}
% head: A. Grothendieck : Dérivateurs, Ch. II, Édition des Universités de Montpellier II et Paris VII
% foot: Version 8 juillet 2001
(which can be interpreted as $\Fl(\underline{\Fl}(\mathcal{M}))$, the set of arrows in $\underline{\Fl}(\mathcal{M})$). We have a subset $\mathcal{D}_0$ of $\mathcal{D}$, consisting of the diagrams for which \texttt{[there]} exists \texttt{[an]} $f$ with two commutative triangles (it is thus the image of $\underbrace{\Fl_3(\mathcal{M})}_{\mathrm{Hom}(\boldsymbol{\Delta}^3,\mathcal{M})}$ in $\mathrm{Hom}(\boldsymbol{\Delta}^1\times\boldsymbol{\Delta}^1,\mathcal{M})$, transpose of $\boldsymbol{\Delta}^1\times\boldsymbol{\Delta}^1\longrightarrow\boldsymbol{\Delta}^3$ given by $(0,0)\longmapsto 0$, $(0,1)\longmapsto 1$, $(1,0)\longmapsto 2$, $(1,1)\longmapsto 3$), and a map
\[
\mathcal{D}_0\subseteq\mathcal{D}\xrightarrow{\ \varphi\ }\Fl(\mathcal{M})\times\Fl(\mathcal{M}),
\]
associating to $(*)$ the pair $(i,j)$ \texttt{[rather} $(i,p)$\texttt{]} (source and target of $(*)$, considered as an arrow in $\underline{\Fl}(\mathcal{M})$ from $i$ to $p$). The relation $\mathcal{R}$ is then deduced from $\varphi$, $\mathcal{D}_0$ by the condition
\msp{page 82}
\[
\Phi\longleftrightarrow\Psi\qquad\Longleftrightarrow\qquad\varphi^{-1}(\Phi\times\Psi)\subseteq\mathcal{D}_0.
\]
But to tell the truth, I realize that it is after all deduced from a relation between elements of $\Fl(\mathcal{M})$, and this relation between $i : A\longrightarrow B$ and $p : X\longrightarrow Y$ can be written as follows:
\begin{itemize}[leftmargin=3em]
\item[$(**)$] for all arrows $A\longrightarrow X$, $B\longrightarrow Y$ making the diagram $(*)$ commutative, the dotted arrow $\varphi$ in $(*)$ exists. % [sic] la flèche $\varphi$ -> la flèche $f$, the dotted arrow in $(*)$ is labelled $f$
\end{itemize}
\end{itemize}

This makes me think that if we have a relation $\mathcal{R}$ between $\Pf(P)$ and $\Pf(Q)$, satisfying $(\mathrm{D}_0)$, hence described by a decreasing map
\[
\Pf(P)\longrightarrow\Pf(Q)
\]
transforming sup into inf, the datum of such a map is equivalent to that of an arbitrary map % [sic] la donné -> la donnée, feminine noun
\[
P\longrightarrow\Pf(Q),
\]
or what amounts to the same thing, of a relation $\mathcal{R}$ on $P\times Q$. Thus this structure must be regarded as ``of the first kind''. The situation changes from the moment when, to describe the relation $\mathcal{R}$ with the desired properties, we find ourselves compelled to replace $\Pf(P)$, $\Pf(Q)$ by suitable subsets - when we look only at certain subsets of $P$, $Q$. This is what we
\msp{page 83}are going to do in the next §. The case I am going to treat there is the only one of its kind that I have encountered - a case of ``formal duality'' which comes in naturally, and which is not ``of the 1\textsuperscript{st} kind''.

\printedpage{51}
% head: A. Grothendieck : Dérivateurs, Ch. II, Édition des Universités de Montpellier II et Paris VII
% foot: Version 8 juillet 2001
\section*{6\quad Transporters, kernels, stabilizers relative to $W$}

Let $\mathcal{M}$ be a category, equipped with a subset
\[
W\subseteq\Fl(\mathcal{M}).
\tag*{(6.1)}
\]
I am thinking of the case of a localizer, but will introduce an assumption on $W$ only as needed. For simplicity, I assume $\mathcal{M}$ stable under fiber products.

I set
\[
\left\{
\begin{array}{rcl}
\mbox{(\footnotemark)}\mathcal{P} & = & \Pf(W)=\{A\in\Pf(\Fl(\mathcal{M}))\mid A\subseteq W\}\\[1ex]
\mathcal{Q} & = & \{L\in\Pf(\Fl(\mathcal{M}))\mid\
\parbox[c]{17em}{$L$ \textit{stable} under base\\
change (by which I also mean,\\
stable under isomorphism)\\
and under composition. $\}$.}
\end{array}
\right.
\tag*{(6.2)}
\]
\footnotetext{cf. correction p. 86.}
I introduce a relation
\[
\mathcal{R}\subseteq\mathcal{P}\times\mathcal{Q}
\]
in the following way. For $A\in\mathcal{P}$, $L\in\mathcal{Q}$, we set
\[
(A,L)\in\mathcal{R}\iff
\begin{array}[t]{l}
\text{For every cartesian diagram}\\
\begin{tikzcd}[ampersand replacement=\&]
X \arrow[d, "f"] \& X' \arrow[l] \arrow[d, "f'"] \\
Y \& Y' \arrow[l, "p"]
\end{tikzcd}\\
\text{with $f\in A$, $p\in L$, we have $f'\in W$}\\
\text{(i.e. $A$ is stable under base change by $p\in L$).}
\end{array}
\tag*{(6.3)}
\]

\msp{page 84}When I order $\mathcal{P},\mathcal{Q}\subseteq\Pf(\Fl(\mathcal{M}))$ by the order relation induced by inclusion of subsets of $\Fl(\mathcal{M})$, we find ordered subsets of $\Pf(\Fl(\mathcal{M}))$, stable under arbitrary unions, and also under arbitrary intersections, with the sole proviso that in the case of $\mathcal{P}$, the family of indices must be $\neq\emptyset$ (since then we find as intersection the whole of $\Fl(\mathcal{M})$, which is not in general contained in $W$). % [sic] la seul réserve -> la seule réserve, feminine agreement

We are going to verify condition $(\mathrm{D}_0)$ on page 69, for $\mathcal{R}$ and the order relations. There are two checks to make.

1$^{o}$) \texttt{[where is ``2$^{o}$''?]} Let us give ourselves $A\in\mathcal{P}$, i.e. $A\subseteq W$, and consider the set $\mathcal{R}(A)$ of $L\in\mathcal{Q}$ such that we have $(A,L)\in\mathcal{R}$. It is clear that this is an open set (and likewise, that ${}^{s}\mathcal{R}(L)$ is an open set of $\mathcal{Q}$ for every $L$ in $\mathcal{Q}$) - i.e. $L\in\mathcal{R}(A)$, $L'\subseteq L$ implies % ?? uncertain: left superscript before R(L) (also on p. 53) -- italic s or upright s % [sic] ouvert de $\mathcal{Q}$ -> ouvert de $\mathcal{P}$, cf. "l'ouvert ${}^{s}\mathcal{R}(L)$ de $\mathcal{P}$" p. 53
\printedpage{52}
% head: A. Grothendieck : Dérivateurs, Ch. II, Édition des Universités de Montpellier II et Paris VII
% foot: Version 8 juillet 2001
$L'\in\mathcal{R}(A)$. It remains to find a greatest element of $\mathcal{R}(A)$. Let us note that if $L\in\mathcal{R}(A)$, then for every $p\in L$, $p : Y'\longrightarrow Y$, the base change
\[
p^{*} : \mathrm{Cat}/Y\longrightarrow\mathrm{Cat}/Y' : X\longmapsto X\times_{Y}Y'
\]
must transform arrows $f\in A$ into arrows $g\in W$, as follows from the fact that $L$ is stable under base change.
\[
\begin{tikzcd}
X_1 \arrow[d, "f"'] & X'_1 \arrow[l, "q_1"'] \arrow[d, "g"] \\
X_2 \arrow[d] & X'_2 \arrow[l, "q_2"'] \arrow[d] \\
Y & Y'. \arrow[l, "p"']
\end{tikzcd}
\]
Thus $L\subseteq\mathbf{L}_W(A)$, where we set
\msp{page 85}
\[
\begin{aligned}
\mathbf{L}_W(A)=\{(p : Y'\longrightarrow Y)\in\Fl(\mathcal{M})\quad\mid\quad & p^{*} : \mathrm{Cat}/Y\longrightarrow\mathrm{Cat}/Y'\ \text{transforms}\\
& \text{arrows of $A$ into arrows \textit{of} $W$.}\}
\end{aligned}
\]
(\footnote{NB say it rather in terms of cartesian diagrams, since we have made no assumption of stability of $W$ under isomorphisms.}).
\[
\begin{tikzcd}
X_1 \arrow[d, "f"'] \arrow[dr, phantom, "\text{cart.}"] & X'_1 \arrow[l] \arrow[d, "f'"] \\
X_2 \arrow[d] \arrow[dr, phantom, "\text{cart.}"] & X'_2 \arrow[l] \arrow[d] \\
Y & Y' \arrow[l, "p"']
\end{tikzcd}
\]
\[
f\in A\implies f'\in W.
\]
It is clear on the other hand that this subset of $\Fl(\mathcal{M})$ is stable under base change, hence $\mathbf{L}_W(A)\in\mathcal{Q}$, and that $(A,\mathbf{L}_W(A))\in\mathcal{R}$. Thus $\mathbf{L}_W(A)$ is a greatest element of $\mathcal{R}(A)$, so we find
\[
\mathcal{R}=\{(A,L)\in\mathcal{P}\times\mathcal{Q}\mid L\subseteq\mathbf{L}_W(A)\}.
\tag*{(6.5)}
\]
\printedpage{53}
% head: A. Grothendieck : Dérivateurs, Ch. II, Édition des Universités de Montpellier II et Paris VII
% foot: Version 8 juillet 2001
\[
\begin{tikzcd}
X_1 \arrow[d, "f"'] \arrow[dr, phantom, "\text{cart.}"] & X'_1 \arrow[l, "q_1"'] \arrow[d, "f'"] \\
X_2 \arrow[d] \arrow[dr, phantom, "\text{cart.}"] & X'_2 \arrow[l, "q_2"'] \arrow[d] \\
Y_1 \arrow[d] \arrow[dr, phantom, "\text{cart.}"] & Y'_1 \arrow[l, "p_1"'] \arrow[d] \\
Y & Y'. \arrow[l, "p"']
\end{tikzcd}
\]
Conversely, let us start with $L\in\mathcal{Q}$, and show that the open set ${}^{s}\mathcal{R}(L)$ of $\mathcal{P}$ has a greatest element. For this, let us set
\[
\begin{aligned}
\mathbf{A}_W(L)=\{(f : X\longrightarrow Y)\in W\subseteq\Fl(\mathcal{M})\quad\mid\quad\forall\qquad
&\begin{tikzcd}[ampersand replacement=\&]
X \arrow[d, "f"'] \arrow[dr, phantom, "\textit{cart.}"] \& X' \arrow[l, "q"'] \arrow[d, "f'"] \\
Y \& Y', \arrow[l, "p"]
\end{tikzcd}\\
&\text{with $p\in L$, we have $f'\in W$}\}.
\end{aligned}
\tag*{(6.6)}
\]
(NB When $L$ contains the identities, as is the case in particular if it is of the form $\mathbf{L}_W(A)$ - then it contains the isomorphisms, at least if $A$ contains them - then $\mathbf{A}_W(L)$ is the set of $f\in\Fl(\mathcal{M})$ satisfying the condition stated after $|$, without having to specify $f\in W$.) We have $\mathbf{A}_W(L)\in\mathcal{P}$, and we verify that $(\mathbf{A}_W(L),L)\in\mathcal{R}$ (using the stability of $L$ under composition), finally it is clear that $(A,L)\in\mathcal{R}$ implies $A\subseteq\mathbf{A}_W(L)$ (using $A\subseteq W$), so $\mathbf{A}_W(L)$ is a greatest element of ${}^{s}\mathcal{R}(L)$.

\msp{page 86}Let us note that if $W$ \textit{is stable under composition} (which we are going to assume from now on), the same is true of $\mathbf{A}_W(L)$ (which again uses the fact that $L$ is stable under base change). Moreover, if $W$ \textit{contains the isomorphisms} (which we shall also assume), $\mathbf{A}_W(L)$ contains them too. As for $\mathbf{L}_W(A)$, it contains the isomorphisms \textit{provided that} $A$ is stable under isomorphisms; when $A$ is stable under composition and contains the isomorphisms, this condition will automatically be satisfied. We insist on this condition, that is why on the side of $A$, in the definition of $\mathcal{P}$, I am going to introduce this additional condition that $A$ be \textit{stable under isomorphisms.} So I set, correcting (6.2),
\[
\left\{
\begin{array}{rcl}
\mathcal{P} & = & \{A\in\Pf(\Fl(\mathcal{M}))\quad\mid\quad\mathrm{Is}(\mathcal{M})\subseteq A\subseteq W,\quad A\text{ stable under composition}\}\\[1ex]
\mathcal{Q} & = & \{L\in\Pf(\Fl(\mathcal{M}))\quad\mid\quad\mathrm{Is}(\mathcal{M})\subseteq L,\quad L\text{ stable under base change}\}.
\end{array}
\right.
\tag*{(6.7)}
\]